% Осень26!A35; нумерация внутри столбца без пустых строк.
% Источник: лекция 7, §5 (строка 146)
\ticket{32}{Лемма о единственности плотности}
\label{ticket:32}

\begin{lemma}
Пусть $f,g$ суммируемы на пространстве с мерой $(X,\mathcal A,\mu)$ и
\[
  \int_A f\,d\mu=\int_A g\,d\mu\qquad(A\in\mathcal A).
\]
Тогда $f=g$ почти везде.
\end{lemma}

\begin{proof}
По \hyperref[ticket:28]{конечности суммируемых функций почти везде} можно изменить $f,g$ на общем измеримом множестве меры нуль так, чтобы они были конечны всюду; интегралы при этом не меняются.
Положим $h=f-g$.
По \hyperref[ticket:36]{линейности интеграла} $h$ суммируема и $\int_A h\,d\mu=0$ для всех $A\in\mathcal A$.

Пусть $A_n^+=\{h\ge1/n\}$ и $A_n^-=\{h\le-1/n\}$.
По монотонности
\[
  0=\int_{A_n^+}h\,d\mu\ge\frac1n\mu(A_n^+),
\]
\[
  0=\int_{A_n^-}(-h)\,d\mu\ge\frac1n\mu(A_n^-).
\]
Следовательно, все $A_n^\pm$ имеют меру нуль, а
\[
  \{h\ne0\}=\bigcup_{n=1}^{\infty}(A_n^+\cup A_n^-)
\]
также имеет меру нуль.
\end{proof}
